\section{Generalized pairs and auxiliary programs}\label{sec:preliminaries}

The relative constructions used in the proof take place on models with a
stronger global factoriality property than the one in
Theorem~\ref{thm:main}.  We first specify that property and the generalized
pair conventions for those constructions.  We then record the precise
forms of the auxiliary results that will be used.

The auxiliary definitions and statements in this section use reduced,
irreducible, second-countable complex analytic spaces as their ambient
spaces.

\begin{definition}[Global strong factoriality]\label{def:strong-factoriality}
A normal compact complex space \(T\) is \emph{globally strongly
\(\Q\)-factorial} if, for every coherent rank-one reflexive sheaf
\(\mathcal F\) on all of \(T\), some positive reflexive power
\[
 \mathcal F^{[m]}=(\mathcal F^{\otimes m})^{**},\qquad m>0,
\]
is an invertible sheaf.
\end{definition}

This condition implies global Weil \(\Q\)-factoriality: apply it to the
divisorial sheaf of each global prime divisor and to the reflexive canonical
sheaf.  It is a condition on sheaves on the whole space; it makes no assertion
that arbitrary analytic open subsets, or the analytic local rings, are
\(\Q\)-factorial.  We use \emph{strong model} to mean a globally strongly
\(\Q\)-factorial model.

A \emph{rational line bundle} on \(T\) is an element of
\(\Pic(T)\otimes_{\Z}\Q\).  An equality of rational line bundles means an
isomorphism of holomorphic line bundles after taking a positive common
integral multiple.  A rank-one reflexive sheaf is \emph{rationally
invertible} if a positive reflexive power is invertible.  On a projective
morphism \(T\to V\), a rational line bundle is \emph{relatively nef} if
its degree is nonnegative on every compact irreducible curve in a fiber,
and is \emph{relatively ample} if a positive integral multiple is relatively
ample.  For the absolute nef data in this paper, nefness means that the
first Chern class lies in the closure of the K\"ahler cone, as in the
introduction; it implies relative nefness for every morphism under
consideration.

The auxiliary constructions are stated for rational generalized
\emph{b-line data}.  On a normal compact space \(T\), these consist of an
effective rational Weil divisor \(B\) of finite support, a projective
modification \(p:W\to T\), and a rational line bundle \(M_W\) on \(W\)
that is nef absolutely, or nef over the specified base of a relative
argument.  It is pulled back on every higher model; we continue to denote
the resulting b-object by \(\M\).  After clearing a denominator, proper
direct image followed by double dual defines its reflexive trace \(M_T\).
The adjoint \(D_T=K_T+B+M_T\) is required to be rationally invertible.
Here \(K_T\) in the b-line formulation is the reflexive canonical sheaf.
On a smooth log resolution projective over \(T\) and carrying the data, the natural
meromorphic comparison over the isomorphism locus defines the crepant
boundary by
\[
 K_W+B_W+M_W=p^*D_T.
\]
The log discrepancy of a global place \(E\) is
\(a(E;T,B+\M)=1-\operatorname{coeff}_E(B_W)\) on a model carrying it.
All statements about places below concern the global places defined in the
introduction.  A global place is exceptional over \(T\) if its center on
\(T\) has codimension at least two.

We abbreviate \emph{generalized klt}, \emph{generalized lc}, and
\emph{generalized dlt} to gklt, glc, and gdlt.  A pair is gklt if all of
its global log discrepancies are positive, and glc if they are all
nonnegative.  It is \emph{generalized terminal} if it is gklt and
\(a(E;T,B+\M)>1\) for every place exceptional over \(T\).  A
\emph{generalized log canonical center} is the center of a global place of
log discrepancy zero.  A glc pair is gdlt if there is a Zariski-open set
\(T^\circ\subset T\) on which \((T^\circ,B|_{T^\circ})\) is simple normal
crossing and the b-data descend, and which contains the general point of
every generalized log canonical center.  The carrier and subsequent log
resolutions may be chosen isomorphic on this open.  The zero-discrepancy
centers there are strata of the coefficient-one boundary.  These are the
conventions of \cite[Convention~9.1 and Definition~9.2]{KahlerAuxiliary}.

We explain the relation between this formulation and the actual divisors in
Theorem~\ref{thm:main}.  A \(\Q\)-Cartier divisor determines a rational
line bundle.  If \(mM_W\) is integral Cartier, then on the normal target
the reflexive sheaf
\[
 \bigl(p_*\mathcal O_W(mM_W)\bigr)^{**}
\]
agrees with the divisorial sheaf \(\mathcal O_T(mp_*M_W)\): the two agree
at every codimension-one point, where the modification is an isomorphism,
and normality gives uniqueness of the reflexive extension.  The same
comparison on a common higher model is compatible with pullback and with
transport across a small map.  The canonical trace is the canonical
sheaf of the target.  These facts are
\cite[Lemma~9.3(i)--(ii)]{KahlerAuxiliary}.  The natural meromorphic
comparisons therefore agree with the comparisons using the compatible
actual canonical representatives in \eqref{eq:discrepancy-definition}.
In every application below the b-line discrepancies are exactly the
discrepancies of that equation, and a crepant b-line identity is the
corresponding equality of actual rational divisors.  A common projective
carrier for any finite part of our birational constructions carries the
pullback of the original analytically nef divisor, which is relatively nef
over each base used in that part.  A single carrier dominating an infinite
sequence is not required.

Here is the elementary step used in the auxiliary programs.  Suppose that
the source \(T\) and the output \(T_1\) are strong normal compact K\"ahler
spaces, and write \(D_T\) and \(D_{T_1}\) for their rationally invertible
adjoints.  An \emph{elementary negative step} is either a projective
bimeromorphic divisorial contraction \(f:T\to T_1\) for which \(-D_T\)
is \(f\)-ample, or a diagram
\[
 T\xrightarrow{\ f\ }Z\xleftarrow{\ f^+\ }T_1
\]
of projective small bimeromorphic morphisms to a normal compact K\"ahler
space with \(-D_T\) \(f\)-ample and \(D_{T_1}\) \(f^+\)-ample.  The
degree functionals of the curves contracted on the negative side span one
nonzero ray when paired with global rational line bundles on \(T\).
In the small case the positive side is the relative Proj of the actual
adjoint algebra: after choosing a positive multiple of \(D_T\) represented
by a holomorphic line bundle \(\mathcal L\), it is
\[
 T_1=\Proj_Z\bigoplus_{m\geq0} f_*\mathcal L^{\otimes m}.
\]
The boundary is pushed forward in the divisorial case and strictly
transformed in the small case, and the same b-data are used on common higher
models.  With the compatible actual canonical representatives, a divisorial
step satisfies \(D_{T_1}=f_*D_T\).  If the step is \emph{over \(V\)}, all its maps commute with the
given maps to \(V\).  In the constructions used below the contraction maps
have connected fibers, and a small step has nonisomorphic morphisms on both
sides; the latter properties are also verified in
Section~\ref{sec:comparisons}.  This is the form of
\cite[Definition~9.6]{KahlerAuxiliary}.  Its one-ray condition concerns
degrees of global line bundles and imposes no condition on the relative
Bott--Chern dimension.  A diagram in Theorem~\ref{thm:main} need not be
elementary.

\begin{proposition}[Low places and projective realization]\label{prop:low-places}
Let \((T,B+\M)\) be gklt rational generalized data on a normal compact
K\"ahler space, with effective boundary of finite support and an absolutely
nef rational datum carried by a projective modification.  Fix a projective
log resolution carrying the datum.  There is \(\varepsilon>0\) such that
every global place satisfies
\[
 a(E;T,B+\M)\geq\varepsilon,
\]
and only finitely many places exceptional over \(T\) satisfy
\[
 a(E;T,B+\M)<1+\varepsilon.
\]
This finite set can be represented simultaneously by prime divisors on a
smooth model projective over \(T\).  In particular the exceptional places
of discrepancy at most one form a finite set with such a realization.
\end{proposition}

This is the rational nef specialization of
\cite[Lemma~4.2(i)]{KahlerAuxiliary}.  The positive number
\(\varepsilon\) belongs to this one pair; uniformity along a sequence will
follow from discrepancy monotonicity.

\begin{proposition}[Exact extraction]\label{prop:extraction}
Let \((T,B+\M)\) be rational gklt data on a strong normal compact
K\"ahler space, with an absolutely nef rational line bundle on a projective
carrier and rationally invertible adjoint.  For any specified finite set
\(\mathcal E\) of places exceptional over \(T\) with
\(a(E;T,B+\M)\leq1\), there is a projective bimeromorphic morphism
\(h:U\to T\) such that \(U\) is a strong normal compact K\"ahler space,
the \(h\)-exceptional prime divisors are exactly the members of
\(\mathcal E\), and the crepant generalized pair on \(U\) is gklt with
effective boundary.  Thus, writing this boundary as \(B_U\),
\[
 K_U+B_U+M_U=h^*(K_T+B+M_T).
\]
\end{proposition}

This is \cite[Proposition~9.11]{KahlerAuxiliary}.  Terminality is not part
of the extraction assertion.  It will follow when \(\mathcal E\) is chosen
to contain every exceptional place of discrepancy at most one.

\begin{proposition}[Relative programs]\label{prop:relative-program}
Let \(\pi:T\to V\) be a projective bimeromorphic morphism of normal compact
K\"ahler spaces.  Suppose that \(T\) is strong and that \((T,B+\M)\) is
rational gdlt, with effective boundary, a projective carrier on which the
rational line datum is nef over \(V\), and a rationally invertible adjoint
\(D_T\).
\begin{enumerate}
\item If \(D_T\) is not relatively nef over \(V\), there is an elementary
negative step over \(V\).  Its output is projective over \(V\), normal
compact K\"ahler and strong, and is gdlt with the transported boundary and
b-data.  The same construction is available after every finite prefix whose
last adjoint is not relatively nef over \(V\).
\item If the initial pair is gklt, there exists a finite, possibly empty,
sequence of elementary negative steps over this one fixed base \(V\) whose
last adjoint is relatively nef over \(V\).  Its models and transported data
have the properties in the first assertion.
\end{enumerate}
\end{proposition}

The two assertions are respectively
\cite[Propositions~9.7 and 9.9]{KahlerAuxiliary}.  The first is a
continuation assertion for gdlt data.  The second asserts the existence of
one finite run for gklt data; it does not assert termination of every choice
of a relative program.

\begin{proposition}[Crepant gdlt modification]\label{prop:gdlt-modification}
Let \((T,B+\M)\) be rational glc data of dimension at most four on a normal
compact K\"ahler space.  Assume that the higher absolutely nef rational
line datum is carried by a projective modification and that
\(K_T+B+M_T\) is rationally invertible.  There is a projective bimeromorphic
morphism \(h:Y\to T\) such that \(Y\) is a strong normal compact K\"ahler
space, \((Y,B_Y+\M)\) is gdlt with \(B_Y\geq0\), and
\[
 K_Y+B_Y+M_Y=h^*(K_T+B+M_T).
\]
Every \(h\)-exceptional prime has coefficient one in \(B_Y\), equivalently
has generalized log discrepancy zero over \((T,B+\M)\).
\end{proposition}

This is \cite[Proposition~11.4]{KahlerAuxiliary}.  The base \(T\) need not
be strong, and the assertion does not require that every zero-discrepancy
place be extracted.

\begin{theorem}[Special termination]\label{thm:special-termination}
Let \(d\leq4\), and let
\[
 (T_0,B_0+\M)\dashrightarrow(T_1,B_1+\M)
 \dashrightarrow(T_2,B_2+\M)\dashrightarrow\cdots
\]
be a sequence of elementary steps for rational gdlt data on strong normal
compact K\"ahler spaces of dimension \(d\).  The boundaries are transported
as in an elementary program, and the higher absolutely nef b-line datum is fixed and is
carried by projective modifications.  There is an index \(i_0\) such that,
for every \(i\geq i_0\), the exceptional loci on both sides of the step are
disjoint from every generalized log canonical center on their respective
sides.  The contraction bases in this sequence may vary.
\end{theorem}

This is \cite[Theorem~11.3]{KahlerAuxiliary}.  We use it for the elementary
sequences on strong models.  Full termination for four-dimensional gdlt
data will be deduced in Section~\ref{sec:completion}.

\paragraph{Scope of the auxiliary inputs.}
The results recalled above use the projective relative arguments in
\cite[Sections~9--11]{KahlerAuxiliary}, together with its low-place
calculation.  Relative continuation is Proposition~9.7 there; the finite
run over a fixed bimeromorphic base in Proposition~9.9 uses the local
finite-generation argument of Lemma~9.8, and exact extraction in
Proposition~9.11 uses that finite run.  The special-termination and gdlt
modification arguments in Section~11 form an induction on dimension,
using the threefold termination theorem of Section~10.  These proofs do
not invoke the canonical bundle formula or the absolute
supporting-contraction induction of Section~6.  Indeed, Section~6 uses
the threefold termination theorem as an input to its good-model
construction; the dependency is in that direction.  Thus the use of
\cite{KahlerAuxiliary} here does not require its absolute contraction or
minimal-model theorems.

The imports from \cite{CanonicalKahler} are likewise restricted to its
common-model construction (Proposition~2.1), analytic negativity
(Lemma~2.2), and terminal smoothness in codimension two (Lemma~2.3).
Its fourfold termination theorem is not an input.  All relative
constructions above retain their stated strong-factoriality hypotheses;
these are not additional hypotheses on the given sequence in
Theorem~\ref{thm:main}.
